\documentclass[11pt]{article}
\usepackage[a4paper,margin=24mm]{geometry}
\usepackage{amsmath,amssymb,hyperref}
\hypersetup{colorlinks=true,urlcolor=blue}
\title{Closing the Gaussian source and selected-period certificate gates}
\author{Independent assistance to the general Gaussian thread}
\date{16 November 2026}
\begin{document}\maketitle

\noindent\textbf{Outcome and scope.} There is a certificate route for pointwise Gaussian all-multiplier safety that needs neither an actual-image kernel nor completeness of a computer algebra integration presentation. Its finite data are an original degree-lowering Euler recurrence, original scalar integration-by-parts witnesses for finitely many markers, and certified original moment jets. The accompanying deliberately slow rational implementation generates these by bounded linear searches and returns UNKNOWN on an exhausted cap. Uncapped termination uses classical holonomic specializability and finite generation; no new priority is claimed for those ingredients or for decidability. This is a complete certificate contract and a small implemented independent route, not closure of all Gaussian normal-form questions, general formal coefficient-oracle questions, HC4, or exceptional Mathieu.

The user asked to help the most promising agent toward full closure. The Gaussian thread had just completed the polynomial-symbol Bernstein backend and was implementing its general Gaussian source/Hom producer. Its owned files and mathematical ledger were not edited here. Messages supplied the lemmas below as they became concrete. Proof-assurance flags remain unchanged.

\section*{1. Universal Taylor duality}
Let $D=K[t]\langle\partial\rangle$, $[\partial,t]=1$, over any characteristic-zero field. For EVERY left $D$-module $N$ there is a natural isomorphism
\[
 \operatorname{Hom}_D(N,K[[t]])\simeq
 \operatorname{Hom}_K(N/tN,K).                                      \tag{1}
\]
Here $N/tN$ is an ordinary vector-space cokernel, not a quotient by a left differential submodule. If $\ell(tN)=0$, the inverse is
\[
 \psi_\ell(h)=\sum_{m\ge0}\ell(\partial^m h)t^m/m!.
\]
Differentiation commutes with this formula. The identity
$\partial^m t=t\partial^m+m\partial^{m-1}$ proves
$\psi_\ell(th)=t\psi_\ell(h)$. Conversely evaluation at zero recovers $\ell$, and the Taylor coefficients of a differential map are exactly the displayed ones. No convergence, regularity, or faithful completion is needed.

Write $H_P$ for the exact source from the retained Gaussian formal-completion proof, with actual period map $\Phi_P$. That proof gives $H_P/tH_P\simeq K$ and one-dimensional formal Hom. If $q:N\twoheadrightarrow H_P$ is a PROVED all-amplitude surjection and $\dim_K N/tN=1$, (1) gives
\[
 P\text{ safe}\quad\Longleftrightarrow\quad
 \operatorname{Hom}_D(N,K[t])\ne0.                                \tag{2}
\]
Indeed every formal map on $N$ is a scalar of $\Phi_P q$; over a field a nonzero scalar is invertible. Surjectivity gives both directions. The extra dimension condition is essential for this unmarked version: $D/D[\partial(\partial-1)]$ has a polynomial quotient but also the selected solution $e^t$, and its time-zero quotient has dimension two.

\section*{2. Elementary exact polynomial Hom on the time line}
Suppose $N=D^r/\sum_j D r_j$ has a finite presentation. For each generator supply a nonzero operator $A_i$ with an explicit LEFT identity
\[
 A_i e_i=\sum_j W_{ij}r_j.                                      \tag{3}
\]
The backend verifies (3) by exact PBW multiplication. These identities bound every coordinate of a polynomial map. For
$A=\sum c_{ab}t^a\partial^b$, put
\[
 s_+=\max(a-b),\qquad
 c_+(z)=\sum_{a-b=s_+}c_{ab}(z)_b .
\]
The polynomial $c_+$ is nonzero: the falling factorials have distinct degrees. A nonzero polynomial solution of degree $d$ must satisfy $c_+(d)=0$, by its highest-degree coefficient. If $d+s_+<0$, every relevant falling factorial is already zero, giving the same necessary condition. Thus the maximum nonnegative integer root is an exact sufficient degree bound; no such root means only the zero polynomial solution.

Use these finite coordinate ansatzes, impose every supplied relation $r_j$, and take a rational nullspace. This computes the FULL polynomial Hom of the supplied presentation. It is not a bounded extrapolation. For a holonomic $D_1$ module each cyclic generator has a nonzero annihilator, since a free $D_1$ cyclic submodule would have larger characteristic dimension. Enumerating finite identities (3) therefore terminates when uncapped. The implementation also generates these witnesses by linear algebra, instead of requiring a general Hom/duality backend.

\href{https://arxiv.org/abs/math/0001064}{Oaku--Takayama--Tsai, Polynomial and rational solutions of holonomic systems} is primary background for polynomial solution algorithms. \href{https://arxiv.org/abs/math/0007139}{Tsai--Walther, Computing homomorphisms between holonomic D-modules} supplies the more general established algorithm. The argument above is the elementary time-line specialization, not a new general Hom theorem.

\section*{3. Original jets remove the extra formal-dimension hypothesis}
There is a stronger marked version of (2). Define
\[
 s_-=\min(a-b),\qquad
 c_-(z)=\sum_{a-b=s_-}c_{ab}(z)_b .
\]
Every nonzero formal solution of $A f=0$ has its initial order among the nonnegative integer roots of $c_-$. This follows from the lowest-degree coefficient, with the same falling-factorial boundary argument as above. Let $J_A$ be the maximum such root, or $-1$ if there are none. Jets through degree $J_A$ inject the entire formal solution space. In the no-root case that space is zero.

Consequently, for a proved actual map $N\to H_P$ and (3), a polynomial map on $N$ agrees with the ACTUAL period map if and only if its coordinate jets through $J_{A_i}$ agree. All those jets are finite original Wick calculations. A rational augmented linear system therefore decides whether the selected full map is polynomial. A failed system retains an explicit left-nullspace functional whose value on original jets is nonzero. With an all-amplitude surjection, polynomiality is equivalent to Gaussian safety. Exact equality of $N$ with $H_P$ and one-dimensional formal Hom are unnecessary.

The implementation rejects the $e^t$ false-positive cover using its original constant and linear coefficients. It also accepts a genuine overpresentation $O\oplus Oe^t\twoheadrightarrow O$, with the second summand mapped to zero: its formal solution space has dimension two, but the selected map is polynomial. Full-amplitude coverage remains indispensable. A cyclic unit period $\Phi_P(1)=1$ alone says nothing about omitted multipliers.

\section*{4. The original Euler recurrence certifies ALL multiplier coverage}
This is the strongest combination with the Gaussian agent's new recurrence work. For standard covariance set
\[
 F=tP(X)-\tfrac12\sum_i X_i^2,\quad
 \Theta=-\sum_i X_i\partial_{X_i}-n .
\]
For any nondegenerate rational symmetric covariance $C$, replace the quadratic term by $-X^TC^{-1}X/2$; the same argument holds for its normalized Wick functional. Split circular pairs are included, allowing nonzero safe rational phases in complex Gaussian coordinates.

Suppose an ORIGINAL polynomial differential identity is supplied:
\[
 [b(\Theta)+R]e^F=0,\qquad
 R=\sum_\nu r_\nu t^{a_\nu}\partial_t^{b_\nu}
                         X^{\alpha_\nu}\partial_X^{\beta_\nu},
 \quad\delta_\nu=|\beta_\nu|-|\alpha_\nu|\ge1.              \tag{4}
\]
The literal identity can be checked after dividing by $e^F$, using polynomial arithmetic in $(t,X)$. No completeness of an integration presentation is invoked. Spatial formal integration by parts has $\Theta^*=\sum_i X_i\partial_{X_i}$. For a homogeneous marker $Q$ of degree $q$, (4) gives
\[
 b(q)\Phi_P(Q)=-\sum_\nu r_\nu t^{a_\nu}\partial_t^{b_\nu}
       \Phi_P\big((-\partial_X)^{\beta_\nu}(X^{\alpha_\nu}Q)\big).       \tag{5}
\]
Each marker on the right is homogeneous of degree $q-\delta_\nu$, or is zero. Let $L$ be the largest nonnegative integer root of $b$. Such a root exists: $b(0)=0$ follows from (5) for $Q=1$ and $\Phi_P(1)(0)=1$. For $q>L$, division by $b(q)$ and induction on $q$ prove that the original monomials $X^\gamma$, $|\gamma|\le L$, generate ALL actual amplitude periods over $D_t$. Thus
\[
 P\text{ safe}\quad\Longleftrightarrow\quad
 \Phi_P(X^\gamma)\in K[t]\text{ for all }|\gamma|\le L.        \tag{6}
\]
This is a global completeness proof, not a finite moment-prefix criterion.

If those finitely many periods have degrees $d_\gamma$, set
\[
 c=\max\left(0,\max_\nu\frac{a_\nu-b_\nu}{\delta_\nu}\right),\quad
 e=\max\left(0,\max_{|\gamma|\le L}(d_\gamma-c|\gamma|)\right).
\]
Equation (5) proves $\deg_t\Phi_P(Q)\le c\deg Q+e$ by the same induction. Hence
\[
 \mathbb E(QP^m)=0\quad\text{whenever }m>c\deg Q+e.           \tag{7}
\]
Zero periods are harmless. This is a sufficient bound, not a least bound. An additional verified pure-null premise can absorb an intercept in the usual way.

\paragraph{Existence and coefficient descent.}
The rank-one polynomial exponential module is holonomic. Partial Fourier transform in $X$ preserves holonomicity. Its standard $V$-filtration specializability along the Fourier-coordinate origin gives a nonzero scalar polynomial $b$ and a relation of the form (4): the Fourier Euler operator becomes $\Theta$, and strictly lower $V$-degree becomes positive original spatial drop. See \href{https://arxiv.org/html/math/9805006v1}{Oaku--Takayama, Algorithms for D-modules}, Section 4 and Algorithm 4.6, followed by the restriction theorem in Section 5. These are classical algebraic statements, with no regularity or analytic convergence requirement. For rational input, existence over an algebraic closure descends: the finite linear identity equations have rational coefficients, and their solution space has a nonzero $b$-projection over the extension if and only if it does over $\mathbb Q$.

The independent enumerator searches finite $b$ and $R$ ansatzes, requiring the spatial-drop condition in advance, and checks the original identity. The agent's adapted-basis implementation is a faster source of the same data; it can use this direct checker or its retained explicit Weyl membership witnesses.

\section*{5. Every low-marker scalar equation can also have an ORIGINAL certificate}
For $Q=X^\gamma$ define
\[
 \nabla_i=\partial_{X_i}-(C^{-1}X)_i+tP_i.
\]
A sufficient scalar-operator certificate is the finite polynomial identity
\[
 A(t,\partial_t+P)Q=\sum_i\nabla_i B_i(t,X).                 \tag{8}
\]
It proves $A\Phi_P(Q)=0$ by original integration by parts. For time-independent $Q$, a PBW monomial $t^a\partial_t^b$ on the left is simply $t^aP^bQ$.

The implemented linear search enumerates coefficients of $A$ and the polynomial $B_i$, and checks (8) literally. Holonomicity of the exact source ensures a nonzero $A$ for each $[Q]$; membership in its defining spatial cokernel supplies polynomial $B_i$. Therefore uncapped enumeration terminates. Apply the infinity and origin arguments of Sections 2--3 to decide this actual scalar period from a finite jet. A NO-MATCH certificate gives a fixed original bad marker, with infinitely many nonzero mixed moments. It is a pure-null counterexample only if the unit period is separately certified to equal one.

Combining (4), (6), and (8) gives a general certificate-only pointwise algorithm over rational nondegenerate Gaussian data. It does not require the full integration relations to be independently reproduced. No general safe normal-form classification follows from this decidability application.

\section*{6. Executable evidence and handoff}
All files are next to this note:
\begin{itemize}
\item \texttt{gaussian\_polynomial\_hom\_sidecheck\_20261116.py}: finite-presentation Hom, generated left witnesses and infinity bounds.
\item \texttt{gaussian\_selected\_jets\_sidecheck\_20261116.py}: origin bounds, actual-map jet matching, explicit mismatch functional, rational symmetric-covariance Wick engine.
\item \texttt{gaussian\_ibp\_annihilator\_sidecheck\_20261116.py}: original polynomial certificates (8) and actual selected-period test.
\item \texttt{gaussian\_euler\_cover\_sidecheck\_20261116.py}: original recurrence (4), finite low-marker decision (6), and numerical bound (7).
\item \texttt{verify\_gaussian\_polynomial\_hom\_sidecheck\_20261116.py}: 118 focused exact controls, including independent direct-exponential recurrence replay and original-polynomial IBP replay.
\end{itemize}
The complete small controls accept $P=0$, reject standard $P=X$ and $P=X^2$, and accept the NONZERO pure safe split phase $P=Z$ with $\mathbb E(ZW)=1$. The last has the certified bound $m>\deg Q$. The quadratic control retains $b(z)=z(z-1)$ and both low degrees; unit-only cyclic coverage is not assumed. The controls also exercise punctual modules, a spurious polynomial quotient, cross-generator relations, origin resonances, false witnesses, and UNKNOWN on capped searches.

The default caps are deliberate. No broad phase/degree/support campaign was run. General uncapped termination is a written argument from the cited classical existence theorems, not an inference from the examples. This package does not assert independent review of every earlier Gaussian theorem, a new novelty claim, a new pure counterexample, or complete closure of the user's broad programme. The practical Gaussian thread remains the owner of integration into its canonical implementation and registry.
\end{document}
